\section{Reproducibilidad}
\label{anexo:reproducibilidad}

Este anexo resume las condiciones necesarias para reproducir el análisis
estadístico, las tablas, las figuras y el informe. La fuente computacional
principal del proyecto es:

\begin{center}
    \texttt{notebook/MCC002\_Grupo5\_Notebook.ipynb}
\end{center}

El notebook utiliza rutas relativas, fija las semillas de los procedimientos
aleatorios y genera automáticamente los resultados almacenados en las
carpetas \texttt{resultados/figuras} y \texttt{resultados/tablas}.

\subsection{Entorno computacional}

El análisis fue desarrollado en Windows 10 con Python 3.12.10. Las
dependencias completas se registran en el archivo
\texttt{requirements.txt}. Las principales versiones empleadas fueron:

\begin{center}
\begin{tabular}{ll}
\toprule
Componente & Versión \\
\midrule
Python & 3.12.10 \\
NumPy & 2.4.6 \\
pandas & 3.0.3 \\
SciPy & 1.17.1 \\
statsmodels & 0.14.6 \\
scikit-learn & 1.9.0 \\
seaborn & 0.13.2 \\
\bottomrule
\end{tabular}
\end{center}

No se requiere una GPU. El análisis puede ejecutarse completamente en CPU.

\subsection{Datos y configuración}

El proyecto utiliza el conjunto
\textit{Statlog German Credit Data}. El notebook busca inicialmente:

\begin{center}
    \texttt{data/german.data}
\end{center}

Si el archivo no se encuentra disponible, se descarga desde el UCI Machine
Learning Repository. Después de la carga se verifica que el conjunto
contenga \(1000\) observaciones, \(20\) predictores y una variable de clase,
sin valores faltantes ni filas completamente duplicadas.

La respuesta se recodifica como:

\[
\text{mal crédito}
=
\begin{cases}
0, & \text{si la clase original es }1,\\
1, & \text{si la clase original es }2.
\end{cases}
\]

La transformación produce \(700\) observaciones de buen riesgo y \(300\)
de mal riesgo.

La semilla global utilizada es:

\[
\texttt{SEED}=2026.
\]

Esta semilla controla el bootstrap, la partición de entrenamiento y prueba,
la validación cruzada y el algoritmo Metropolis--Hastings.

\subsection{Ejecución del notebook}

Desde la raíz del repositorio, el entorno puede reconstruirse en PowerShell
mediante:

\begin{verbatim}
python -m venv .venv
.\.venv\Scripts\Activate.ps1
python -m pip install --upgrade pip
python -m pip install -r requirements.txt
\end{verbatim}

El notebook puede ejecutarse desde Visual Studio Code mediante
\textit{Run All} o automáticamente con:

\begin{verbatim}
python -m jupyter nbconvert `
  --to notebook `
  --execute notebook/MCC002_Grupo5_Notebook.ipynb `
  --output MCC002_Grupo5_Notebook_ejecutado.ipynb `
  --output-dir notebook `
  --ExecutePreprocessor.timeout=0
\end{verbatim}

La ejecución debe finalizar sin excepciones y sin modificar manualmente las
salidas entre celdas.

\subsection{Resultados esperados}

La ejecución completa debe generar:

\begin{itemize}
    \item once figuras PNG en \texttt{resultados/figuras};
    \item cuatro tablas CSV en \texttt{resultados/tablas};
    \item las estimaciones, pruebas y métricas utilizadas en el informe.
\end{itemize}

Los principales valores de control son:

\begin{center}
\begin{tabular}{p{7.2cm}p{3.6cm}}
\toprule
Control & Resultado esperado \\
\midrule
Dimensiones del dataset & \(1000\times21\) \\
Valores faltantes & \(0\) \\
Buenos créditos & \(700\) \\
Malos créditos & \(300\) \\
Tasa global de mal crédito & \(0.300\) \\
IC de Wilson del \(95\,\%\) & \([0.2724,\;0.3291]\) \\
AUC en prueba & \(0.7740\) \\
AUC media en validación cruzada & \(0.7685\) \\
Costo con umbral \(0.50\) & \(270\) \\
Costo con umbral \(1/6\) & \(158\) \\
Aceptación de Metropolis--Hastings & \(25.68\,\%\) \\
Media posterior de duración & \(0.04084\) \\
\bottomrule
\end{tabular}
\end{center}

Pueden presentarse pequeñas diferencias numéricas cuando se utilizan
versiones distintas de las bibliotecas, pero las conclusiones y los valores
redondeados deben mantenerse estables.

\subsection{Compilación del informe}

El informe se compila desde la carpeta \texttt{informe} mediante:

\begin{verbatim}
.\limpiar.ps1
.\compilar.ps1
\end{verbatim}

El documento resultante debe generarse en:

\begin{center}
    \texttt{informe/build/main.pdf}
\end{center}

Después de la compilación debe verificarse:

\begin{verbatim}
$LASTEXITCODE
\end{verbatim}

El valor esperado es \(0\). También debe comprobarse que no existan citas
sin resolver, referencias mostradas como \texttt{??} ni errores fatales en
\texttt{build/main.log}.

\subsection{Criterio de reproducibilidad}

El proyecto se considera reproducible porque utiliza rutas relativas,
documenta las dependencias, fija las semillas, descarga o valida los datos,
genera automáticamente sus artefactos y permite reconstruir el informe
mediante scripts de compilación.

La reproducibilidad garantiza que los cálculos puedan repetirse bajo los
mismos datos y supuestos. No implica que los resultados sean directamente
generalizables a otras instituciones, periodos o poblaciones crediticias.